\section{Simultaneous control of all local roof averages}
\label{sec:v91-maximal}
The exceptional-source estimate admits a maximal form. Its conclusion
is simultaneous over all roof averaging radii, but the exceptional
roof set need not be empty. We keep this distinction explicit.

\begin{lemma}[A local maximal inequality in the critical norm]
\label{lem:v91-maximal}
Let $q>1$ and extend $g\ge0$ by zero outside a bounded interval $K$.
Define
\[
 \mathcal Mg(t)=\sup_{I\ni t}\frac1{|I|}\int_I g(u)\,du,
\]
where the supremum is over bounded intervals. If
$\mathcal N_q(g;K)\le A$, then
\begin{equation}\label{eq:v91-maximal-level}
       |\{\mathcal Mg>\lambda\}|\le3(A/\lambda)^q,
                         \qquad \lambda>0.
\end{equation}
For every bounded interval $J$,
$\mathcal N_q(\mathcal Mg;J)\le3^{1/q}q(q-1)^{-1}A$.
\end{lemma}
\begin{proof}
For a compact subset of $\{\mathcal Mg>\lambda\}$ choose a finite
interval cover with averages exceeding $\lambda$. Selecting intervals
in decreasing order of length gives a disjoint subcollection
$I_1,\ldots,I_N$ whose triple dilates cover that compact set.
For $V=\bigcup I_i$, the norm definition and the zero extension give
\[
 \lambda|V|<\int_Vg\le A|V|^{1-1/q}.
\]
Hence $|V|<(A/\lambda)^q$ and the compact set has measure at most
$3(A/\lambda)^q$. Inner regularity proves
\eqref{eq:v91-maximal-level}. For $E\subset J$, integrate the
bound $|E\cap\{\mathcal Mg>\lambda\}|\le
\min(|E|,3(A/\lambda)^q)$ in $\lambda$. Splitting at
$3^{1/q}A|E|^{-1/q}$ gives the asserted norm bound. Rational
interval endpoints suffice by local integrability, so all functions
and exceptional sets can be chosen measurable.
\end{proof}

\begin{theorem}[All-radius estimate on the original source]
\label{thm:v91-maximal-source}
Fix $M,L>0$. For each interval $J$ of length $L$, let
$K=J+(-1,1)$ and put
\[
 G_{m,a,\delta}=m^2U_{a,\delta}
                          \one_{\mathcal C(M)\cap K}.
\]
For every fixed $\lambda>0$,
\begin{equation}\label{eq:v91-maximal-source}
 \limsup_m\sup_{R,n,k}\sup_{|J|=L}
 |\{t\in J:\mathcal M G_{m,a,\delta}(t)>\lambda\}|
                  \le C_{M,L}\lambda^{-Q}
                                      (a^\alpha+\delta^\gamma).
\end{equation}
Outside this set, every interval average contained in $K$ of every
$m^2R_{a',\delta'}^\varepsilon\one_{\mathcal C(M)}$, with
$a'\le a$, $\delta'\le\delta$, is at most $\lambda$ when the
interval contains the starting roof. The same bound holds for the
source values at their Lebesgue points. In the cofinal choice the
right side of \eqref{eq:v91-maximal-source} is
$C_{M,L}\lambda^{-Q}h^{1/128}$.
\end{theorem}
\begin{proof}
Apply Lemma~\ref{lem:v91-maximal} to the actual finite-count source
$G_{m,a,\delta}$ and then take the specified suprema and collision
limsup. Theorem~\ref{thm:v91-exceptional} bounds its norm by
$C_{M,L}(a^\alpha+\delta^\gamma)^{1/Q}$ in that limit.
Raising to the $Q$th power proves \eqref{eq:v91-maximal-source}.
All smaller weighted selectors are dominated by this one unweighted
union using the common coarea kernels. Their interval averages are
therefore bounded simultaneously. The assertion on point values
uses Lebesgue differentiation for each source; no common null set
for all radii $R$ is asserted. The cofinal exponents were computed
in the preceding theorem.
\end{proof}

\subsection{The pointwise obligation that remains}
The organizing assertion is still \eqref{eq:v76-target}. Its
essential-height criterion is the vanishing, in the outer
reconstruction limit, of the original weighted critical-or-short-
residence source charge. Theorems~\ref{thm:v91-raw} and
\ref{thm:v91-exceptional} remove the logarithmic loss at the finite
critical exponent. Theorem~\ref{thm:v91-maximal-source} estimates
one exceptional roof set for all local averaging scales. None
asserts that this set is empty or that its essential height is small.

For example, $\one_{(0,\delta)}$ has essential height one and
critical norm $\delta^{1/Q}$. Its maximal superlevel sets shrink,
and its $W^{s,1}$ norm tends to zero for $s<1$, without any decline
of its essential height. This is a logical test of the topology,
not a proposed billiard counterexample. Thus the present critical
convergence is not identified with a collision-uniform spike-width
lower bound. The latter, or another genuinely pointwise physical
estimate, is still needed for the original scalar theorem.

All actual labels and the original source normalization remain in
these statements. Point-conditioned path ratios and likelihoods
remain downstream of scalar pointwise closure and their retained
positive-arithmetic hypotheses; arithmetic zero classes are never
divided by. The historical integrated-oscillation proofs are retained
in full, not replaced by the maximal statement.
